\section{Экспериментальная оценка}
\label{sec:Chapter5} \index{Chapter5}

Для оценки работы разработанного прототипа была проведена серия экспериментов, в ходе которых разработанный форк Trivy с поддержкой eStargz (в дальнейшем~--- \textit{trivy-estargz}) сравнивался с тремя другими сканерами: оригинальным Trivy, Grype (в связке с Syft) и Clair. В данном разделе подробно описаны методология эксперимента, разработанный для его проведения инструментарий и полученные результаты.

\subsection{Методология бенчмарка}

Целью экспериментальной оценки являлось получение количественных показателей эффективности предлагаемого подхода в условиях, приближённых к реальному использованию. Для этого требовалось решить ряд методологических вопросов:

\begin{itemize}
\item какие образы использовать в качестве тестового набора, чтобы он отражал реальный профиль нагрузки;
\item какие именно метрики собирать в ходе сканирования;
\item как обеспечить воспроизводимость результатов;
\item как подтвердить корректность работы прототипа~--- что результаты сканирования совпадают с результатами обычного Trivy.
\end{itemize}

\subsubsection{Тестовый набор образов}

В качестве модельного профиля нагрузки рассматривается \textit{managed container registry} облачного провайдера, через который проходят образы, регулярно создаваемые автоматизированными CI/CD-конвейерами. Такой сценарий обладает следующими характерными свойствами:

\begin{itemize}
\item образы строятся на основе разнообразных базовых образов~--- Ubuntu, Alpine, Fedora, специализированных образах для конкретных языков (Python, Node.js, Go);
\item образы делятся на относительно стабильную часть (базовый Linux + системные пакеты + языковые зависимости) и часто изменяющуюся часть (приложение);
\item один и тот же образ накатывается на registry в различных версиях с одинаковым набором зависимостей, но разным содержимым приложения~--- что моделирует обновления приложения при неизменных зависимостях.
\end{itemize}

Для воспроизводимости результатов был разработан скрипт \texttt{generate\_images.py}, генерирующий тестовый набор образов по описанному выше шаблону. Скрипт принимает на вход следующие параметры:

\begin{itemize}
\item \textit{количество производных образов на базовый}~--- сколько различных образов на каждом базовом дистрибутиве сгенерировать (в наших экспериментах~--- 5);
\item \textit{количество версий каждого производного}~--- сколько <<коммитов>> сэмулировать для каждого производного образа (в наших экспериментах~--- 3);
\item \textit{случайный seed}~--- для детерминированной генерации.
\end{itemize}

С такими параметрами скрипт генерирует $6 \times 5 \times 3 = 90$ образов, покрывающих следующие типы:

\begin{itemize}
\item \texttt{ubuntu:22.04} + apt-пакеты;
\item \texttt{alpine:3.18} + apk-пакеты;
\item \texttt{fedora:39} + dnf-пакеты;
\item \texttt{python:3.11-slim} + apt-пакеты + pip-пакеты;
\item \texttt{node:20-slim} + apt-пакеты + npm-пакеты;
\item \texttt{golang:1.23} с дальнейшей упаковкой Go-бинарного файла в \texttt{scratch}-образ.
\end{itemize}

Для каждого производного образа набор устанавливаемых пакетов генерируется случайным образом из заранее заданных пулов реальных популярных пакетов (например, для apt~--- curl, vim, openssh, libssl-dev и т.п.; для pip~--- requests, flask, numpy и т.п.). Это обеспечивает реалистичный набор зависимостей в каждом образе.

Случай Go-образа заслуживает отдельного внимания: чтобы получить реалистичный образ с метаинформацией о модулях, вшитой в бинарный файл, используется multi-stage сборка~--- сначала исходный файл собирается во временный builder-контейнер на базе \texttt{golang:1.23}, затем готовый бинарный файл переносится в \texttt{scratch}-образ. Это даёт ровно тот сценарий, который для Trivy является самым неудобным с точки зрения частичной загрузки: метаинформация о Go-модулях находится внутри самого бинарного файла, который занимает значительную часть образа.

Для каждого образа также генерируются разные версии приложения (различные исходные файлы), что обеспечивает разные верхние слои при общем составе зависимостей.

\subsubsection{Метрики}

В рамках бенчмарка с каждого сканера для каждого образа собираются следующие метрики:

\begin{itemize}
\item \textit{wall\_time}~--- полное время сканирования образа (от запуска до получения отчета);
\item \textit{index\_time}~--- время инвентаризации (там, где это разделение явно выражено~--- для Grype это время работы Syft, для Clair~--- время работы indexer);
\item \textit{match\_time}~--- время сопоставления с базой уязвимостей;
\item \textit{vulns}~--- количество найденных уязвимостей по уровням критичности (Critical, High, Medium, Low, Other);
\item \textit{layers\_fetched}~--- список дайджестов слоев, скачанных из registry в ходе сканирования;
\item \textit{layers\_cached}~--- список дайджестов слоев, взятых из локального кеша;
\item для trivy-estargz дополнительно~--- \textit{bytes\_fetched}: для каждого слоя~--- реальный объем скачанных из registry данных, как замеренный счетчиком в \texttt{httpRangeReader}.
\end{itemize}

Сбор метрик с разных сканеров требует индивидуального подхода. У Trivy и Grype подходящие сведения извлекаются из их вывода в режиме \texttt{--debug}/JSON. У Clair всё несколько сложнее: метрики о фактически скачанных слоях получаются из логов контейнера indexer~--- они парсятся регулярными выражениями по характерным строкам логов.

\subsubsection{Условия запуска и воспроизводимость}

Все сканеры запускались на одной и той же машине, в одной и той же сети относительно registry. Перед запуском каждого сканера полностью очищался его локальный кеш (\texttt{cold} режим), чтобы все слои образа гарантированно скачивались из registry, и измерения отражали \textit{наихудший} с точки зрения скорости сценарий~--- т.е. сценарий, в котором эффект от использования eStargz проявляется наиболее ярко.

Базы уязвимостей перед началом эксперимента предварительно загружались для каждого сканера, и в режиме \texttt{cold} не очищались~--- очистка касалась только кеша слоев. Это разделение обусловлено тем, что в реальных сценариях база уязвимостей загружается один раз и используется при сканировании множества образов, поэтому её повторная загрузка при каждом эксперименте не отражает реальный профиль работы.

\subsection{Разработка тестового фреймворка}

Для проведения описанных экспериментов был разработан фреймворк, состоящий из нескольких компонент-скриптов:

\begin{itemize}
\item \texttt{generate\_images.py}~--- генерация и загрузка тестового набора в registry, описанная выше;
\item \texttt{convert\_estargz.py}~--- конвертация всех образов из тестового набора в формат eStargz и их загрузка в registry под отдельным тегом-суффиксом (\texttt{-estargz});
\item \texttt{run\_benchmark.py}~--- основной runner бенчмарка, выполняющий сканирование всех образов всеми указанными сканерами и собирающий метрики в JSON-файл;
\item \texttt{compare\_trivy.py}~--- скрипт верификации, сравнивающий результаты сканирования trivy и trivy-estargz по каждой паре <<образ, severity>>;
\item \texttt{analyze.py}~--- скрипт построения графиков по собранным результатам.
\end{itemize}

\subsubsection{Архитектура фреймворка}

Фреймворк построен по модели <<генерация $\rightarrow$ конвертация $\rightarrow$ сканирование $\rightarrow$ верификация $\rightarrow$ анализ>>. Каждый из шагов формирует свой артефакт, и каждый следующий принимает на вход артефакт предыдущего:

\begin{itemize}
\item генератор формирует \texttt{manifest.json} со списком сгенерированных образов и их метаинформацией (базовый образ, тип производного, версия, состав слоев, размер);
\item конвертатор использует \texttt{manifest.json} для определения списка образов и для каждого создает eStargz-вариант;
\item runner также принимает на вход \texttt{manifest.json}, обходит каждый образ каждым из заданных сканеров и записывает собранные метрики в \texttt{results-<scanner>.json};
\item верификатор и аналитик работают со всеми собранными результатами и формируют итоговые отчеты~--- текстовые расхождения и графики соответственно.
\end{itemize}

Все скрипты написаны на Python и не имеют сложных зависимостей сверх стандартной библиотеки и нескольких популярных пакетов (\texttt{requests}, \texttt{matplotlib}, \texttt{numpy}). Это упрощает запуск фреймворка на различных машинах и в различных средах. Все рандомизированные шаги (выбор пакетов, версий) контролируются параметром \texttt{seed} и тем самым обеспечивают полную воспроизводимость генерации тестового набора.

\subsubsection{Развертывание окружения}

В составе фреймворка предполагается работа со следующими внешними компонентами:

\begin{itemize}
\item \textit{container registry}~--- в качестве registry в экспериментах использовался Harbor, развернутый на отдельной машине в той же сети. Harbor выбран как наиболее распространённое open-source решение для корпоративного использования, поддерживающее как OCI-совместимое хранение, так и интеграцию со сканерами;
\item \textit{Clair}~--- развертывается отдельно в виде набора контейнеров (indexer, matcher, postgres) через \texttt{docker-compose};
\item \textit{Trivy, Grype, Syft}~--- устанавливаются локально как обычные CLI-утилиты;
\item \textit{trivy-estargz}~--- собирается из исходного кода разработанного форка и используется как отдельный бинарный файл, чтобы не пересекаться с обычным Trivy.
\end{itemize}

\subsection{Обзор результатов и сравнение}

После прохождения всего пайплайна были собраны метрики по сканированию каждого из 90 образов всеми четырьмя сканерами. Графики ниже представляют усреднённые результаты по типу базового образа.

\subsubsection{Производительность сканирования}

На рис.~\ref{fig:perf} приведено среднее время сканирования (\texttt{wall\_time}) для каждого сканера и каждого типа базового образа.

\begin{figure}[H]
    \centering
    \includegraphics[width=0.95\textwidth]{chart-performance.png}
    \caption{Среднее время сканирования по типу базового образа}
    \label{fig:perf}
\end{figure}

Можно заметить следующие наблюдения:

\begin{itemize}
\item \texttt{trivy-estargz} существенно быстрее обычного \texttt{trivy} на всех типах образов, кроме \texttt{golang}~--- разница в среднем составляет 4--10X;
\item \texttt{clair} оказывается сопоставим по скорости с \texttt{trivy-estargz}, что объясняется его архитектурой~--- он использует внутренний кеш слоев и параллельную обработку;
\item \texttt{grype} в большинстве случаев медленнее обычного \texttt{trivy}, что объясняется отсутствием в нём кеша слоев и тем, что для каждого образа полностью производится инвентаризация;
\item на \texttt{golang}-образах разница между \texttt{trivy} и \texttt{trivy-estargz} существенно меньше~--- это объясняется тем, что для Go-образов в формате \texttt{scratch} сканеру нужен весь бинарный файл (как обсуждалось ранее, метаинформация о Go-модулях вшита в бинарный файл и не выделяется в отдельный манифест).
\end{itemize}

Стоит отметить, что наблюдаемое ускорение \texttt{trivy-estargz} относительно \texttt{trivy} носит не линейный, а скачкообразный характер~--- при определенном пороге размера слоя время сканирования выходит на плато (несколько секунд на образ), определяемое прежде всего временем работы анализаторов и матчера на содержимом, а не сетевыми задержками. Этот эффект особенно заметен на \texttt{alpine}-образах, где обычный Trivy всё равно скачивает значительные объемы данных, в то время как \texttt{trivy-estargz} ограничивается единицами мегабайт.

\subsubsection{Объём скачиваемых данных}

На рис.~\ref{fig:size} приведены три величины, усреднённые по типу базового образа: \textit{estargz fetched data size}~--- средний объем данных, реально скачанных trivy-estargz из registry; \textit{estargz image size}~--- размер eStargz-варианта образа в registry; \textit{original image size}~--- размер исходного (не eStargz) варианта образа в registry.

\begin{figure}[H]
    \centering
    \includegraphics[width=0.95\textwidth]{chart-image-size.png}
    \caption{Размеры образов и объёмы скачиваемых данных по типу базового образа}
    \label{fig:size}
\end{figure}

Видно, что для большинства типов образов trivy-estargz скачивает существенно меньше данных, чем составляет размер самого образа. Также можно заметить, что размер eStargz-варианта образа на 3--8\% больше исходного~--- это плата за пофайловое сжатие. Этот overhead вполне допустим и в значительной степени окупается экономией трафика при сканировании.

Те же данные в другом виде~--- в виде отношений~--- приведены на рис.~\ref{fig:ratio}. На графике показано отношение объёма скачанных trivy-estargz данных к двум размерам: \textit{fetched / estargz image size} и \textit{fetched / original image size}.

\begin{figure}[H]
    \centering
    \includegraphics[width=0.95\textwidth]{chart-fetch-ratio.png}
    \caption{Доля скачанных данных от размера образа по типу базового образа}
    \label{fig:ratio}
\end{figure}

В среднем при использовании trivy-estargz скачивается в районе 5--15\% от объёма исходного образа, что согласуется с теоретическими оценками, приведенными в \hyperref[sec:Chapter3]{разделе 4}. Стоит отдельно прокомментировать характерные различия между типами образов:

\begin{itemize}
\item для \texttt{golang}-образов отношение близко к 100\%, поскольку в текущей реализации анализатора Go-модулей требуется весь бинарный файл целиком. Это ожидаемое поведение, и в данной ситуации применение ленивой загрузки не даёт выигрыша;
\item для образов \texttt{alpine} и \texttt{fedora} отношение выше, чем для \texttt{ubuntu}~--- это связано с тем, что в этих образах содержимое \texttt{/usr/lib} и других подобных директорий, не используемых при сканировании, занимает меньшую часть от общего размера образа, поэтому процент скачиваемых для анализа данных оказывается выше;
\item для \texttt{node}-образов отношение оказывается ниже, чем для \texttt{python}, поскольку существенную часть node-образов занимают \texttt{node\_modules}, а сканер анализирует только манифесты \texttt{package.json} и \texttt{package-lock.json}, что показывает наибольшую эффективность ленивой загрузки.
\end{itemize}

На рис.~\ref{fig:scatter} приведена scatter-диаграмма, показывающая зависимость времени сканирования от объёма скачанных данных для каждого образа в отдельности (только для trivy-estargz). Видно сильную корреляцию между этими величинами, причём наклон линии тренда меньше для \texttt{golang}-образов (т.е. больше времени тратится на анализ содержимого, нежели на скачивание).

\begin{figure}[H]
    \centering
    \includegraphics[width=0.85\textwidth]{chart-time-to-image-size.png}
    \caption{Зависимость времени сканирования от объёма скачанных данных (trivy-estargz)}
    \label{fig:scatter}
\end{figure}

\subsubsection{Корректность результатов}

Помимо производительности, не менее важной является корректность результатов сканирования. Для её верификации использовался скрипт \texttt{compare\_trivy.py}, попарно сравнивающий результаты сканирования обычного Trivy и trivy-estargz для каждого образа.

Сравнение производилось по количеству уязвимостей в разбивке по уровням критичности (Critical, High, Medium, Low, Other). По результатам верификации все 90 образов показали \textit{полное} совпадение количества уязвимостей между двумя сканерами по всем уровням критичности. Это подтверждает, что предложенный подход не нарушает корректности работы сканера: использование ленивой загрузки даёт ровно те же результаты, что и обычное скачивание слоев.

\subsection{Обсуждение результатов}

Полученные результаты позволяют сделать следующие основные выводы:

\begin{itemize}
\item предложенный подход обеспечивает кратное (в среднем 4--10X) ускорение сканирования относительно обычного Trivy в режиме работы с registry;
\item объём скачиваемых из registry данных снижается до 5--15\% от размера образа в большинстве случаев;
\item корректность результатов сканирования сохраняется в полной мере~--- предложенный подход не вводит каких-либо приближений или потерь информации;
\item за это приходится платить незначительным (3--8\%) увеличением размера образа в формате eStargz по сравнению с обычным форматом.
\end{itemize}

При этом эффективность подхода зависит от типа сканируемого образа. Наибольший выигрыш наблюдается для образов с большим объемом языковых пакетов (особенно \texttt{node\_modules}), наименьший~--- для образов, в которых анализ требует обработки больших бинарных файлов (Go в формате scratch).

Стоит отдельно отметить, что в сравнении с другими сканерами trivy-estargz демонстрирует производительность, сопоставимую с Clair, при этом сохраняя архитектурную простоту Trivy~--- сканер по-прежнему остаётся stateless CLI-утилитой, не требующей дополнительной инфраструктуры (PostgreSQL, фоновых процессов и т.п.). Это делает предложенный подход особенно привлекательным для встраивания в существующие CI/CD-конвейеры и serverless-окружения.

\subsubsection{Ограничения экспериментальной оценки}

При интерпретации полученных результатов следует учитывать ряд ограничений:

\begin{itemize}
\item тестовый набор содержит синтетически сгенерированные образы, моделирующие типичный профиль нагрузки, но не покрывающие весь спектр реальных сценариев. В частности, в нём отсутствуют образы со специализированным содержимым (например, ML-модели, большие наборы данных, образы Windows-контейнеров и т.п.);
\item эксперимент проводился на одной паре <<сканер--registry>> в локальной сети с малыми задержками. Результаты при работе через интернет с большими задержками могут отличаться: с одной стороны, рост задержек должен сильнее влиять на trivy-estargz, поскольку он использует больше HTTP-запросов; с другой стороны, экономия суммарного объема данных перевешивает рост задержек, особенно для крупных образов;
\item базы уязвимостей разных сканеров отличаются по составу, что приводит к расхождениям в количестве найденных уязвимостей между сканерами. Это нормально и не считается ошибкой: сравнение производительности проводилось без оглядки на абсолютные значения;
\item бенчмарк проводился на сравнительно небольшом наборе образов (90 единиц). В реальной эксплуатации количество образов в registry на порядки больше, что в общем случае увеличивает влияние кеша слоев. Все запуски в данной работе проводились с пустым кешем (\textit{cold}), что соответствует наихудшему сценарию.
\end{itemize}

Несмотря на перечисленные ограничения, проведённая экспериментальная оценка позволяет сделать достаточно уверенный вывод о применимости предложенного подхода и его потенциале для практического внедрения.

\newpage
